% !TEX root = ../main.tex
\begin{ouverture}[lunette]{18}{Lentilles minces et images}{Partie 4 · Ondes et signaux}
{Lentille convergente\enspace·\enspace Relation de conjugaison\enspace·\enspace Grandissement\enspace·\enspace Image réelle ou virtuelle\enspace·\enspace Mise au point}
\objectif{Exploiter les relations de conjugaison et de grandissement pour déterminer la position et la taille de l'image d'un objet-plan réel.}
\objectif{Déterminer les caractéristiques d'une image : réelle ou virtuelle, droite ou renversée, agrandie ou réduite.}
\objectif{Estimer la distance focale d'une lentille ; tester la relation de conjugaison.}
\objectif{Expliquer une mise au point par modification de la distance focale ou de la géométrie du montage (œil, appareil photo).}
\prerequis{Lentille mince convergente, foyers, distance focale, image réelle, grandissement (classe de seconde).}
\prerequis{Théorème de Thalès ; grandeurs algébriques sur un axe orienté.}
\end{ouverture}
\essentiel
\essSection{1}{Lentille mince convergente}
\begin{definition}[Foyers, distance focale, vergence]
Une lentille mince convergente, de centre optique O, fait converger en son \textbf{foyer image} $F'$ les rayons parallèles à l'axe optique.
Le \textbf{foyer objet} $F$ est le symétrique de $F'$ par rapport à O. La \textbf{distance focale} vaut $f'=\alg{OF'}>0$ et la \textbf{vergence} $C=\dfrac{1}{f'}$ s'exprime
en dioptries ($\delta$, avec $f'$ en m) : une lentille de \qty{20}{cm} de focale a une vergence de \num{5,0}\,$\delta$.
\end{definition}
\begin{propriete}[Trois rayons particuliers]
Un rayon passant par O n'est pas dévié ; un rayon parallèle à l'axe ressort en passant par $F'$ ; un rayon passant par $F$ ressort
parallèle à l'axe. Deux de ces rayons, issus d'un point B de l'objet, suffisent à construire son image B'.
\end{propriete}
\schema{18-lentille}
\essSection{2}{Relations de conjugaison et de grandissement}
\begin{loi}[Relations de conjugaison et de grandissement]
On oriente l'axe optique dans le sens de propagation de la lumière, avec l'origine en O. Les grandeurs algébriques, surmontées
d'une barre, sont positives dans le sens de l'axe. L'image A'B' d'un objet AB perpendiculaire à l'axe vérifie :
\[\frac{1}{\alg{OA'}}-\frac{1}{\alg{OA}}=\frac{1}{f'}\qquad\text{et}\qquad\gamma=\frac{\alg{A'B'}}{\alg{AB}}=\frac{\alg{OA'}}{\alg{OA}}.\]
Exemple : avec $f'=\qty{10}{cm}$ et $\alg{OA}=\qty{-30}{cm}$, on trouve $\alg{OA'}=\qty{15}{cm}$ et $\gamma=-0{,}5$.
\end{loi}
\begin{propriete}[Caractéristiques de l'image]
Si $\alg{OA'}>0$, l'image est \textbf{réelle} : on peut la recueillir sur un écran. Si $\alg{OA'}<0$, elle est \textbf{virtuelle}, du même côté que l'objet.
Si $\gamma<0$, l'image est \textbf{renversée} ; si $\gamma>0$, elle est \textbf{droite}. Elle est agrandie si $|\gamma|>1$. Un objet réel placé entre $F$ et O donne une
image virtuelle, droite et agrandie : c'est la \textbf{loupe}.
\end{propriete}
\par\penalty0
\schema{18-loupe}
\begin{python}
fp = 10                                   # cm : distance focale
for OA in [-50, -30, -20, -15, -12, -5]:  # cm : position de l'objet
    OAp = 1 / (1 / fp + 1 / OA)           # relation de conjugaison
    print(OA, 'cm :', round(OAp, 1), 'cm ; gamma =', round(OAp / OA, 2))
\end{python}
\begin{console}
-50 cm : 12.5 cm ; gamma = -0.25
-30 cm : 15.0 cm ; gamma = -0.5
-20 cm : 20.0 cm ; gamma = -1.0
-15 cm : 30.0 cm ; gamma = -2.0
-12 cm : 60.0 cm ; gamma = -5.0
-5 cm : -10.0 cm ; gamma = 2.0
\end{console}
\noindent\capacite{Capacité numérique} Quand l'objet s'approche du foyer, l'image réelle s'éloigne et grandit ; entre le foyer et la lentille, elle
devient virtuelle et droite.
\essSection{3}{Mise au point : l'œil et l'appareil photo}
\begin{propriete}[Former une image nette]
Pour qu'une image soit nette, elle doit se former sur le capteur. Dans un \textbf{appareil photo}, la distance lentille-capteur est
réglable : la mise au point consiste à déplacer l'objectif. Dans l'\textbf{œil}, la distance cristallin-rétine est fixe : c'est la vergence du
cristallin qui augmente quand on regarde de près (\textbf{accommodation}). Certains appareils utilisent des lentilles à focale variable.
\end{propriete}
\schema{18-oeil}
\begin{attention}
Respecter les signes : pour un objet réel placé avant la lentille, $\alg{OA}<0$. Utiliser les mêmes unités pour $\alg{OA}$, $\alg{OA'}$ et $f'$ ; la vergence
se calcule avec $f'$ en mètres.
\end{attention}
\begin{experience}[Tester la relation de conjugaison]
Sur un banc d'optique, on forme sur un écran l'image nette d'un objet lumineux pour plusieurs positions de l'objet. On trace
$\dfrac{1}{\alg{OA'}}$ en fonction de $\dfrac{1}{\alg{OA}}$ : on obtient une droite de pente 1, d'ordonnée à l'origine $\dfrac{1}{f'}$.
\end{experience}
\begin{rediger}
Écrire la relation littérale, isoler l'inconnue, puis remplacer en respectant les signes et les unités. Conclure sur la nature de
l'image (réelle ou virtuelle, droite ou renversée, agrandie ou réduite) à partir des signes de $\alg{OA'}$ et de $\gamma$.
\end{rediger}
\exercices

\rubrique{Découverte}
\begin{blocLivre}{ex-18-1}
\begin{exercice}[1]{Distance focale et vergence}
\begin{questions}
\item Rappeler la définition du foyer image d'une lentille convergente.
\item Calculer la vergence d'une lentille de distance focale \qty{20}{cm}.
\item Calculer la distance focale d'une lentille de vergence \num{8,0}\,$\delta$, puis d'une lentille de \num{50}\,$\delta$.
\item Laquelle de ces lentilles est la plus convergente ? Justifier.
\end{questions}
\end{exercice}
\end{blocLivre}
\begin{blocLivre}{ex-18-2}
\begin{exercice}[1]{Construire une image}
\schema{18-ex02}
\begin{questions}
\item Lire la distance focale de la lentille et la position de l'objet.
\item Construire l'image A'B' de l'objet AB à l'aide de trois rayons particuliers.
\item Mesurer $\alg{OA'}$ et $\alg{A'B'}$. En déduire le grandissement.
\item Décrire l'image : réelle ou virtuelle, droite ou renversée, agrandie ou réduite.
\end{questions}
\end{exercice}
\end{blocLivre}
\begin{blocLivre}{ex-18-3}
\begin{exercice}[1]{Calculer une position et une taille d'image}
Un objet AB de \qty{2,0}{cm} est placé à \qty{30}{cm} devant une lentille convergente de distance focale \qty{10}{cm}.
\begin{questions}
\item Que vaut $\alg{OA}$ ? Justifier son signe.
\item Calculer $\alg{OA'}$ à l'aide de la relation de conjugaison.
\item Calculer le grandissement, puis la taille de l'image.
\item Calculer la position de l'image si l'objet est placé à \qty{15}{cm} de la lentille.
\end{questions}
\end{exercice}
\end{blocLivre}
\begin{blocLivre}{ex-18-4}
\begin{exercice}[1]{La loupe}
On observe une lettre de \qty{2,0}{mm} avec une loupe de distance focale \qty{5,0}{cm}, placée à \qty{3,0}{cm} au-dessus de la lettre.
\begin{questions}
\item Calculer la position de l'image.
\item L'image est-elle réelle ou virtuelle ? Peut-on la recueillir sur un écran ?
\item Calculer le grandissement et la taille de l'image de la lettre.
\item Construire l'image à l'échelle 1/2.
\end{questions}
\end{exercice}
\end{blocLivre}
\begin{blocLivre}{ex-18-5}
\begin{exercice}[1]{Projeter une image}
Avec une lentille, on obtient l'image nette d'une diapositive, placée à \qty{25}{cm} de la lentille, sur un écran situé à \qty{1,00}{m} de la lentille.
\begin{questions}
\item Que valent $\alg{OA}$ et $\alg{OA'}$, en mètres ?
\item Calculer la vergence de la lentille, puis sa distance focale.
\item Calculer le grandissement. L'image est-elle droite ou renversée ?
\item Pourquoi place-t-on les diapositives à l'envers dans un projecteur ?
\end{questions}
\end{exercice}
\end{blocLivre}
\rubrique{Entraînement}
\begin{blocLivre}{ex-18-6}
\begin{exercice}[2]{La relation de conjugaison en Python}
On reprend le programme de l'essentiel.
\begin{questions}
\item \capacite{Capacité numérique}Retrouver la ligne \code{OAp = ...} à partir de la relation de conjugaison.
\item Pour quelle position de l'objet l'image a-t-elle la même taille que l'objet ? Où se trouve-t-elle alors ?
\item Que se passe-t-il quand l'objet s'approche du foyer objet ?
\item Interpréter la dernière ligne affichée.
\end{questions}
\end{exercice}
\end{blocLivre}
\begin{blocLivre}{ex-18-7}
\begin{exercice}[2]{L'appareil photo}
Un appareil photo a un objectif de distance focale \qty{50}{mm} et un capteur de \qty{24}{mm} sur \qty{36}{mm}. On photographie une personne de
\qty{1,80}{m} située à \qty{3,0}{m}.
\begin{questions}
\item Pour un objet très éloigné, où se forme l'image ? Où place-t-on alors le capteur ?
\item Calculer la distance objectif-capteur pour photographier la personne. De combien a-t-on déplacé l'objectif ?
\item Calculer la taille de l'image de la personne sur le capteur. Tient-elle dans la hauteur de \qty{24}{mm} ? Et dans la largeur ?
\item Calculer la distance objectif-capteur pour un objet situé à \qty{1,0}{m}.
\end{questions}
\end{exercice}
\end{blocLivre}
\begin{blocLivre}{ex-18-8}
\begin{exercice}[2]{L'accommodation de l'œil}
On modélise l'œil par une lentille (le cristallin) placée à \qty{17}{mm} d'un écran (la rétine).
\begin{questions}
\item Quand l'œil regarde un objet très éloigné, au repos, quelle est la distance focale du cristallin ? Sa vergence ?
\item Calculer la distance focale et la vergence nécessaires pour voir nettement un objet situé à \qty{25}{cm}.
\item De combien de dioptries l'œil doit-il accommoder ?
\item Calculer la taille, sur la rétine, de l'image d'une lettre de \qty{5}{mm} lue à \qty{25}{cm}.
\end{questions}
\end{exercice}
\end{blocLivre}
\begin{blocLivre}{ex-18-9}
\begin{exercice}[2]{Estimer une distance focale}
Pour estimer la distance focale d'une lentille, on forme sur une feuille l'image d'une fenêtre éloignée : l'image est nette quand la
feuille est à \qty{15}{cm} de la lentille.
\begin{questions}
\item Justifier qu'on obtient ainsi une estimation de la distance focale.
\item On place ensuite un objet à \qty{50}{cm} de la lentille. Où former son image ?
\item Si la fenêtre est en réalité à \qty{5,0}{m}, de combien l'image est-elle décalée par rapport au foyer ?
\item Cette méthode est-elle précise ? Proposer une meilleure méthode.
\end{questions}
\end{exercice}
\end{blocLivre}
\begin{blocLivre}{ex-18-10}
\begin{exercice}[2]{Tester la relation de conjugaison}
Sur un banc d'optique, avec une lentille de focale inconnue, on a mesuré :
\par\smallskip\noindent\hfill\begin{tabular}{lcccccc}
\toprule
$\alg{OA}$ (cm) & $-15$ & $-20$ & $-25$ & $-30$ & $-40$ & $-50$ \\
$\alg{OA'}$ (cm) & 30,2 & 19,8 & 16,6 & 15,1 & 13,4 & 12,4 \\
\bottomrule
\end{tabular}\hfill\null\par\smallskip
\begin{questions}
\item Calculer $\dfrac{1}{\alg{OA'}}-\dfrac{1}{\alg{OA}}$ pour chaque mesure. Que constate-t-on ?
\item Tracer $\dfrac{1}{\alg{OA'}}$ en fonction de $\dfrac{1}{\alg{OA}}$. Quelle courbe obtient-on ?
\item En déduire la distance focale et la vergence de la lentille.
\item Citer une source d'erreur de cette expérience.
\end{questions}
\end{exercice}
\end{blocLivre}
\begin{blocLivre}{ex-18-11}
\begin{exercice}[2]{Démontrer la relation de grandissement}
\schema{18-ex11}
\begin{questions}
\item Pourquoi les droites (AB) et (A'B') sont-elles parallèles ?
\item Appliquer le théorème de Thalès dans les triangles OAB et OA'B'.
\item En tenant compte des signes, en déduire que $\gamma=\dfrac{\alg{A'B'}}{\alg{AB}}=\dfrac{\alg{OA'}}{\alg{OA}}$.
\item Pourquoi une image réelle d'un objet réel est-elle toujours renversée ?
\end{questions}
\end{exercice}
\end{blocLivre}
\rubrique{Objectif DS}
\begin{blocLivre}{ex-18-12}
\begin{exercice}[3]{Le vidéoprojecteur}
\begin{docu}{1}{Le principe}
Dans un vidéoprojecteur, une petite matrice de \qty{14}{mm} de large, fortement éclairée, forme l'objet. Un objectif de distance focale
\qty{20}{mm} en donne une image agrandie sur un écran.
\end{docu}
\begin{docu}{2}{Installation}
L'écran est d'abord placé à \qty{5,0}{m} de l'objectif, puis à \qty{3,0}{m}.
\end{docu}
\begin{questions}
\item Calculer la distance entre la matrice et l'objectif pour un écran à \qty{5,0}{m}.
\item Calculer le grandissement et la largeur de l'image.
\item \initiative L'écran est rapproché à \qty{3,0}{m}. Comment faire la mise au point ? De combien faut-il déplacer l'objectif ?
\item Calculer la largeur de l'image à \qty{3,0}{m}. Comment obtenir une grande image dans une petite salle ?
\end{questions}
\end{exercice}
\end{blocLivre}
\begin{blocLivre}{ex-18-13}
\begin{exercice}[3]{L'appareil photo d'un smartphone}
\begin{docu}{1}{Le module photo}
L'objectif d'un smartphone a une distance focale de \qty{4,0}{mm}. Un petit moteur peut l'éloigner du capteur d'au plus \qty{0,2}{mm} par rapport à
sa position pour un objet à l'infini.
\end{docu}
\begin{questions}
\item Où se trouve le capteur quand on photographie un paysage lointain ?
\item De combien faut-il déplacer l'objectif pour photographier un objet situé à \qty{10}{cm} ?
\item \initiative Calculer la distance minimale de mise au point du smartphone.
\item Calculer le grandissement à cette distance. Pourquoi les smartphones ajoutent-ils parfois un objectif « macro » ?
\end{questions}
\end{exercice}
\end{blocLivre}
\begin{blocLivre}{ex-18-14}
\begin{exercice}[3]{Observer un timbre à la loupe}
\begin{docu}{1}{Observer confortablement}
Un œil normal voit nettement sans fatigue les objets situés au-delà d'environ \qty{25}{cm}. Un philatéliste observe un timbre de \qty{2,5}{cm} de haut
avec une loupe de distance focale \qty{8,0}{cm}, tenue à \qty{6,0}{cm} du timbre.
\end{docu}
\begin{questions}
\item Calculer la position de l'image du timbre donnée par la loupe.
\item Calculer le grandissement et la taille de l'image.
\item L'œil étant placé contre la loupe, l'observation est-elle confortable ? Justifier.
\item \initiative Que se passerait-il si le timbre était placé à plus de \qty{8,0}{cm} de la loupe ?
\end{questions}
\end{exercice}
\end{blocLivre}
\begin{blocLivre}{ex-18-15}
\begin{exercice}[3]{La lentille liquide}
\begin{docu}{1}{Une lentille à focale variable}
Certaines caméras utilisent une lentille liquide : une goutte dont la courbure varie avec une tension électrique. La distance entre la
lentille et le capteur reste fixe, égale à \qty{5,0}{mm}, comme la distance cristallin-rétine dans l'œil.
\end{docu}
\begin{questions}
\item Quelle distance focale faut-il pour photographier un objet très éloigné ? Calculer la vergence correspondante.
\item Calculer la distance focale et la vergence nécessaires pour un objet situé à \qty{20}{cm}.
\item \initiative Quelle variation de vergence faut-il pour passer de l'infini à \qty{10}{cm} ?
\item Comparer ce mode de mise au point à celui d'un appareil photo classique et à celui de l'œil.
\end{questions}
\end{exercice}
\end{blocLivre}
